% META: {"id": "TCOMPL-ECH-ME-005", "chapitre": "TCOMPL-ECHANTILLONNAGE", "type_objet": "methode", "capacites_codes": ["C5"], "methodes": ["M5"], "mode_creation": "ex_nihilo", "sources_inspiration": [], "status": "needs_review"}
% BEGIN-VERIFY
% import random
% random.seed(42)
% ech = [sum(1 for _ in range(100) if random.random() < 0.3) for _ in range(200)]
% moyenne_des_frequences = sum(ech) / len(ech) / 100
% assert abs(moyenne_des_frequences - 0.3) < 0.02
% END-VERIFY
\begin{fichemethode}{M5}{Simuler des echantillons en Python et etudier la serie des moyennes}

\quandUtiliser{%
  Utiliser cette methode lorsque l'enonce demande de simuler en Python un
  echantillon de taille $n$ d'une variable aleatoire, puis de calculer et
  d'observer la serie des moyennes de $N$ echantillons.
}

\pasApas{%
  \begin{enumerate}
    \item Simuler UNE observation : \verb|random.random() < p| renvoie
      \verb|True| (succes) avec probabilite $p$.
    \item Simuler UN echantillon de taille $n$ : compter les succes sur $n$ tirages
      dans une boucle ou une comprehension.
    \item Repeter $N$ fois pour obtenir $N$ echantillons, et stocker chaque
      moyenne (ou frequence) dans une liste.
    \item Calculer la moyenne de la serie des moyennes et son etendue ; comparer
      a la valeur theorique attendue ($p$ pour une frequence de succes).
    \item Observer l'effet de $n$ et de $N$ : les moyennes se resserrent autour de
      la valeur theorique quand $n$ augmente.
  \end{enumerate}
}

\exempleRedige{%
  Simuler $N = 200$ echantillons de taille $n = 100$ d'une Bernoulli de parametre
  $p = 0{,}3$ et calculer la moyenne des frequences observees.

  \medskip
\begin{verbatim}
import random
random.seed(42)                 # simulation reproductible
frequences = []
for _ in range(200):            # N echantillons
    succes = sum(1 for _ in range(100)
                 if random.random() < 0.3)
    frequences.append(succes / 100)
moyenne = sum(frequences) / len(frequences)
\end{verbatim}

  \commentaireMarge{La moyenne des frequences est proche de $p = 0{,}3$.}%
  On obtient une moyenne d'environ $0{,}30$ : la serie des moyennes fluctue
  autour de la valeur theorique $p$.
}

\pieges{%
  \begin{itemize}
    \item \textbf{Piege 1.} Ecrire \verb|random.random() < p| avec $p$ en
      pourcentage (30 au lieu de 0,3).
    \item \textbf{Piege 2.} Confondre $n$ (taille d'un echantillon) et $N$
      (nombre d'echantillons).
    \item \textbf{Piege 3.} Oublier de diviser par $n$ pour passer du nombre de
      succes a la frequence.
  \end{itemize}
}

\verifier{%
  Controler que chaque frequence est dans $[0\,;\,1]$, que la moyenne des
  frequences est proche de $p$, et relancer la simulation avec une autre graine
  pour verifier la stabilite.
}

\sEntrainer{\refExos{M5}}

\end{fichemethode}
